\chapter{Planning and Multi-Step Reasoning}
\label{ch:planning}

\epigraph{``Plans are worthless, but planning is essential.''}{Dwight D. Eisenhower}

\learninggoal{
	\item Distinguish between implicit planning (emergent in the reasoning trace) and explicit planning (structured decomposition).
	\item Implement task decomposition, hierarchical planning, and plan verification.
	\item Understand when and why planning fails, and how to recover.
	\item Compare planning approaches along dimensions of flexibility, reliability, and cost.
}

\section{Why Planning Is Hard}

Planning requires the agent to reason about future states it has not yet observed. This introduces several failure modes not present in reactive loops:

\begin{itemize}
	\item \textbf{Commitment:} choosing a plan early may commit the agent to a suboptimal path.
	\item \textbf{Incomplete information:} the agent may not know what tools or data will be available later.
	\item \textbf{State change:} the environment may change between planning and execution.
	\item \textbf{Decomposition granularity:} too coarse and steps are unmanageable; too fine and overhead dominates.
\end{itemize}

\begin{definition}[Planning]
	\textbf{Planning} is the process of generating a sequence of actions (a plan) that transforms the current state into a goal state, before executing the first action.
\end{definition}

\section{Planning Approaches}

\subsection{Implicit Planning (Emergent)}

The simplest form of planning: the agent reasons step-by-step without an explicit plan structure. This is what ReAct and CoT produce:

\begin{lstlisting}[caption=Implicit planning in ReAct]
  Thought: I need to find the population of Tokyo and its area,
           then compute density.
  Action: search("Tokyo population 2025")
  Observation: 14.1 million
  Thought: Now I need the area.
  Action: search("Tokyo area")
  Observation: 2194 sq km
  Thought: Now I compute: 14.1e6 / 2194
  Action: calculator("14.1e6 / 2194")
  Observation: 6426
  Answer: 6426 people per sq km.
\end{lstlisting}

The plan emerges from the reasoning trace. No separate planning step exists. This is flexible but can produce inconsistent plans (e.g., gathering data in the wrong order, missing prerequisites).

\subsection{Explicit Planning (Decompose Then Execute)}

The agent first generates a complete plan, then executes it step by step:

\begin{lstlisting}[caption=Explicit plan-and-execute]
  # Phase 1: Plan
  Thought: This task requires:
    1. Research: find population and area of Tokyo
    2. Compute: calculate density
    3. Format: present results

  Plan:
    step_1: search("Tokyo population 2025")
    step_2: search("Tokyo area")
    step_3: calculator("population / area")
    step_4: format_answer()

  # Phase 2: Execute
  for step in plan:
      result = execute(step)
      verify(result)
      if failed: replan()
\end{lstlisting}

\subsection{Hierarchical Planning}

Complex tasks decompose into sub-tasks, each potentially planned recursively:

\begin{algorithm}[Hierarchical Planning]
	\begin{enumerate}
		\item Decompose the goal $G$ into sub-goals $\{g_1, \dots, g_n\}$.
		\item For each sub-goal $g_i$:
		      \begin{enumerate}
			      \item If $g_i$ is primitive (directly executable), create an action $a_i$.
			      \item Otherwise, recursively plan for $g_i$.
		      \end{enumerate}
		\item Verify that the plan tree covers all necessary steps.
		\item Execute leaves first, compose results upward.
	\end{enumerate}
\end{algorithm}

This mirrors Hierarchical Task Networks (HTN) from classical planning, but the decomposition is performed by an \llm\ rather than a hand-coded domain model.

\section{Plan Verification}

Plans can be wrong. Verification catches errors before execution:

\begin{lstlisting}[caption=Plan verification prompt]
  system: You are a plan verifier. Check if the following plan
          will achieve the goal. Identify:
          1. Missing steps
          2. Incorrect ordering
          3. Unavailable information
          4. Resource conflicts

  Goal: {goal}
  Plan: {plan}

  Questions:
  - Does each step have the inputs it needs from prior steps?
  - Are there any circular dependencies?
  - Is anything missing?
  - Will the final output meet the goal?
\end{lstlisting}

\section{Failure Recovery}

When a plan fails during execution, the agent must recover:

\begin{longtable}{p{3.5cm}p{4cm}p{4cm}p{3cm}}
	\toprule
	\textbf{Failure mode} & \textbf{Symptom}             & \textbf{Recovery}              & \textbf{Cost}    \\
	\midrule
	Missing information   & Tool returns nothing         & Revise plan, add research step & 1--3 extra calls \\
	Tool error            & Tool returns error           & Retry or find alternative tool & 1 call           \\
	Wrong assumption      & Observation contradicts plan & Replan from current state      & Full re-plan     \\
	Context overflow      & Token limit exceeded         & Summarise, continue            & 1 extra call     \\
	\bottomrule
\end{longtable}

\begin{principle}[Fail Fast, Replan Cheap]
	Verify each step's result immediately after execution. A failed step detected early costs $n$ extra calls. A failed step detected after 10 more steps costs $10n$ extra calls.
\end{principle}

\section{Plan Representation}

Plans can be represented in various formats:

\begin{longtable}{p{3cm}p{4cm}p{4cm}p{3.5cm}}
	\toprule
	\textbf{Format}  & \textbf{Structure}                           & \textbf{Pros}                         & \textbf{Cons}             \\
	\midrule
	Natural language & Free text                                    & Flexible, no parse failures           & Hard to verify, update    \\
	Structured JSON  & \texttt{\{steps: [\{action, args, deps\}]\}} & Machine-parsable, dependency tracking & Rigid schema              \\
	Code             & Executable pseudocode                        & Precise, verifiable by execution      & Overkill for simple tasks \\
	Graph            & DAG of actions                               & Optimal for parallel tasks            & Complex to maintain       \\
	\bottomrule
\end{longtable}

\begin{example}[Structured plan in JSON]
	\begin{lstlisting}
  {
    "goal": "Research and compare Q1 revenue of three companies",
    "steps": [
      {
        "id": 1,
        "action": "search",
        "args": {"q": "AAPL Q1 2025 revenue"},
        "depends_on": [],
        "verify": "result contains a revenue number"
      },
      {
        "id": 2,
        "action": "search",
        "args": {"q": "GOOGL Q1 2025 revenue"},
        "depends_on": [],
        "verify": "result contains a revenue number"
      },
      {
        "id": 3,
        "action": "search",
        "args": {"q": "MSFT Q1 2025 revenue"},
        "depends_on": [],
        "verify": "result contains a revenue number"
      },
      {
        "id": 4,
        "action": "compare_and_format",
        "args": {
          "companies": ["AAPL", "GOOGL", "MSFT"],
          "data_source": "$1, $2, $3"
        },
        "depends_on": [1, 2, 3]
      }
    ]
  }
  \end{lstlisting}
\end{example}

\section{When Planning Fails}

Planning is not always the right approach:

\begin{itemize}
	\item \textbf{Highly dynamic environments:} by the time the plan is formulated, the state has changed. Reactive approaches work better.
	\item \textbf{Unknown unknowns:} the agent cannot plan for things it doesn't know exist. Exploratory approaches (search, trial-and-error) are more robust.
	\item \textbf{Overhead:} planning costs compute. For simple tasks, a direct ReAct loop is cheaper and equally effective.
\end{itemize}

\begin{theorem}[Planning Overhead]
	Explicit planning adds $O(P)$ \llm\ calls where $P$ is the number of plan steps or decomposition depth. The overhead is justified when the task requires more steps than the \llm\ can reliably track in a single trajectory (typically $> 5$--$10$ steps).
\end{theorem}

\section{Exercises}

\begin{exercise}
	Compare implicit (ReAct) and explicit (plan-then-execute) planning on a task with 8 sequential steps. Measure: total \llm\ calls, success rate, and time to completion.
\end{exercise}

\begin{exercise}
	Design a plan verification prompt. Test it on 5 good plans and 5 bad plans (missing steps, wrong order, circular dependencies). What is your verifier's accuracy?
\end{exercise}

\begin{exercise}
	Implement a hierarchical planner for a travel booking task: find flights, compare prices, check visa requirements, book hotel. Show the decomposition tree and the leaf-level actions.
\end{exercise}

\begin{exercise}
	When should you \emph{not} use explicit planning? Design a decision rule that selects between implicit and explicit planning based on task characteristics.
\end{exercise}
